% Full answers PDF: MCQ answers + worked explanations per option, then official MS clips.
% Compiler: XeLaTeX.
\documentclass[12pt,a4paper]{article}
\usepackage[margin=18mm,headheight=14pt]{geometry}
\usepackage{fontspec}
\usepackage[version=4]{mhchem}
\usepackage{graphicx}
\usepackage{xcolor}
\usepackage{booktabs}
\usepackage{fancyhdr}
\usepackage{enumitem}
\usepackage{needspace}

\setmainfont{Helvetica Neue}

\pagestyle{fancy}
\fancyhf{}
\fancyhead[L]{\small 3.6 Intermolecular forces --- answers}
\fancyhead[R]{\small CIE 9701}
\fancyfoot[C]{\thepage}
\renewcommand{\headrulewidth}{0.3pt}

\setlength{\parindent}{0pt}
\setlength{\parskip}{0.35em}

\newcommand{\msblock}[2]{%
  \par\vspace{0.8em}%
  \noindent\begin{minipage}{\textwidth}%
  {\normalsize\bfseries #1}\par\vspace{0.25em}%
  \noindent\includegraphics[width=\textwidth,keepaspectratio]{#2}\par
  \end{minipage}\par
}

\newcommand{\mcq}[4]{%
  \par\needspace{5\baselineskip}\vspace{0.7em}%
  {\normalsize\bfseries #1.\quad #2}%
  \hfill{\footnotesize\color{gray}#3}\par\vspace{0.15em}%
  {\small #4}\par
}

\newcommand{\why}[2]{\par\needspace{2.2\baselineskip}\vspace{0.1em}{\small\textbf{#1} #2}\par}

\begin{document}

\begin{center}
{\large\bfseries 3.6 Intermolecular forces homework --- answers}\\[0.2em]
{\small CIE 9701 AS Chemistry \quad·\quad 8 multiple-choice + 10 written marks}
\end{center}

{\small\bfseries Do not issue this file to students.}
\hrule
\vspace{0.4em}

\section*{Section A --- Multiple choice}

{\normalsize\bfseries Quick answers}

\vspace{0.3em}
{\small
\begin{tabular}{@{}ll@{\hspace{2.6em}}ll@{\hspace{2.6em}}ll@{}}
\toprule
Q & Ans. & Q & Ans. & Q & Ans.\\
\midrule
A1 & \textbf{C} & A4 & \textbf{B} & A7 & \textbf{B}\\
A2 & \textbf{C} & A5 & \textbf{B} & A8 & \textbf{D}\\
A3 & \textbf{B} & A6 & \textbf{A} & & \\
\bottomrule
\end{tabular}}

\vspace{0.6em}
{\normalsize\bfseries Worked explanations}

\mcq{A1}{C}{2021 FM Paper 12 Q5}{%
Temporary dipole--induced dipole forces (instantaneous dipole--induced dipole)
are the \emph{only} intermolecular force in a substance whose molecules have
\emph{no permanent dipole} and \emph{no H-bonding}. Octane, \ce{C8H18}, is a
non-polar hydrocarbon, so this is true for octane.}

\why{A:}{hydrogen chloride is a polar molecule, so it also has permanent
dipole--permanent dipole attractions.}

\why{B:}{methanol, \ce{CH3OH}, has an O--H group, so it shows hydrogen bonding.}

\why{D:}{water has extensive hydrogen bonding.}

\mcq{A2}{C}{2021 MJ Paper 13 Q5}{%
Substance 1 has permanent dipole--dipole attractions $\rightarrow$ a polar
molecule with no H-bonding: \ce{CH2Cl2} (polar C--Cl bonds; no O--H or N--H).
Substance 2 has only weak forces between molecules $\rightarrow$ a non-polar
molecule: \ce{CH4}.
Substance 3 has hydrogen bonding $\rightarrow$ a molecule with an O--H group:
\ce{CH3OH}.
So substance 1 = \ce{CH2Cl2}, substance 2 = \ce{CH4}, substance 3 = \ce{CH3OH}.}

\why{A / B / D:}{mis-match at least one substance --- e.g.\ \ce{Al2Cl6} is
non-polar (it would give only weak forces, not permanent dipole--dipole), and
\ce{H2O} in the ``substance 3'' column would leave \ce{CH3OH} unplaced.}

\mcq{A3}{B}{2024 MJ Paper 11 Q8}{%
Down Group 16 the hydrides \ce{H2S}, \ce{H2Se}, \ce{H2Te} show a steady increase
in boiling point because the molecules get larger with more electrons, so the van
der Waals' forces increase. Water breaks the trend because the O--H bond allows
\emph{strong hydrogen bonding} between water molecules, which is much stronger
than the van der Waals' forces in the other hydrides.}

\why{A:}{shielding/nuclear charge does not explain a \emph{higher} boiling point
anomaly of this size.}

\why{C:}{the covalent O--H bonds are not broken during boiling; only
intermolecular forces are overcome.}

\why{D:}{smaller molecules have \emph{weaker}, not stronger, van der Waals'
forces.}

\mcq{A4}{B}{2024 MJ Paper 12 Q8}{%
Identify the strongest intermolecular force in each compound:
1 \ce{(CH3)3CH} (2-methylpropane) --- non-polar, only van der Waals' forces
(lowest bp, $-$12\,\textdegree C).
4 \ce{CH3CH2CH2CH3} (butane) --- also non-polar, slightly larger so slightly
stronger van der Waals' forces ($-$0.5\,\textdegree C).
3 \ce{CH3CH2CH2SH} (propane-1-thiol) --- weak permanent dipole--dipole
(67\,\textdegree C).
2 \ce{CH3CH2CH2OH} (propan-1-ol) --- hydrogen bonding (97\,\textdegree C).
Increasing boiling point: $1 < 4 < 3 < 2$.}

\why{A / C / D:}{each reverses part of the order --- e.g.\ placing the alcohol
below the thiol, or the thiol below butane.}

\mcq{A5}{B}{2021 ON Paper 12 Q6}{%
Compare the overall dipole moments: \ce{CO2} is linear and symmetrical, so its
bond dipoles cancel and the overall dipole is \emph{zero}. The other three all
have a permanent dipole, and \ce{H2C=O} (methanal) has the largest --- the very
polar \ce{C=O} group with no other polar bond to oppose it
($\mu \approx 2.3$ D, larger than \ce{CH2Cl2} $\approx 1.6$ D and \ce{COCl2}
$\approx 0.9$ D).}

\why{A:}{\ce{CH2Cl2} is polar, but its dipole is smaller than that of \ce{H2C=O}.}

\why{C:}{\ce{COCl2} is polar, but the two C--Cl dipoles partly oppose the C=O
dipole, so its overall dipole is the smaller of the polar molecules here.}

\why{D:}{\ce{CO2} is symmetrical, so it has \emph{no} overall dipole.}

\mcq{A6}{A}{2021 ON Paper 11 Q33}{%
All three statements are correct.
(1) Ethylene glycol molecules insert into and disrupt the extensive
\emph{hydrogen-bond network} of ice, so it melts below 0\,\textdegree C.
(2) Ethylene glycol has two O--H groups, so its molecules hydrogen-bond to each
other.
(3) It also hydrogen-bonds to water molecules, so it dissolves in the meltwater.
Statements 1, 2 and 3 are all correct.}

\why{B / C / D:}{each rejects at least one correct statement. Statement 1 is the
key one: de-icing works precisely because the glycol disrupts the H-bond network
of ice.}

\mcq{A7}{B}{2021 MJ Paper 12 Q3}{%
Each compound has four carbon atoms:
X = butan-1-ol (\ce{CH3CH2CH2CH2OH}) --- hydrogen bonding.
Y = ethoxyethane (\ce{CH3CH2OCH2CH3}) --- permanent dipole--dipole (an ether, no
O--H, so no H-bonding).
Z = sodium butanoate (\ce{CH3CH2CH2COONa}) --- ionic bonding.
The stronger the force broken on boiling, the higher the boiling point:
Y (pd--pd) $<$ X (H-bond) $<$ Z (ionic). So lowest = Y, highest = Z.}

\why{A / C / D:}{these rows put the compounds in the wrong order --- e.g.\
placing the ionic compound Z below the alcohol, or the ether above the alcohol.}

\mcq{A8}{D}{2025 MJ Paper 12 Q10}{%
(1) \textbf{False} --- both $\Delta H_b$ (vaporisation) and $\Delta H_m$
(melting) are \emph{endothermic}, so they are \emph{positive}, not negative.
(2) \textbf{True} --- vaporisation breaks all the intermolecular forces, while
melting only loosens the lattice, so $\Delta H_b > \Delta H_m$.
(3) \textbf{True} --- ice and water are both held together by the same type of
force, hydrogen bonding.
Correct statements: 2 and 3 $\rightarrow$ D.}

\why{A / B / C:}{all include statement 1, which is false because the enthalpy
changes of melting and boiling are positive.}

\newpage

\section*{Section B --- Mark scheme clips}

\msblock{B1  [1]}{cie-9701-2023-mj-p22-q1b-iii-ms.png}

\msblock{B2  [2]}{cie-9701-2023-mj-p22-q1b-iv-ms.png}

\msblock{B3  [2]}{cie-9701-2021-mj-p22-q2c-iii-ms.png}

\msblock{B4  [3]}{cie-9701-2025-mj-p22-q1d-i-ms.png}

\msblock{B5  [2]}{cie-9701-2021-on-p22-q1a-ii-ms.png}

\end{document}
